\section{总结与延伸}

\subsection{核心要点回顾}

\begin{enumerate}
    \item \textbf{扩散语言模型}打破了 AR 模型从左到右的约束，能够并行生成文本
    \item \textbf{MDLM} 通过简洁的连续时间 ELBO，使离散扩散的训练和 AR 模型一样简单，性能在大规模上可匹配甚至超越 AR
    \item \textbf{KV 缓存}是 MDLM 实际部署的最大瓶颈——双向注意力导致激活值无法冻结
    \item \textbf{SoLM} 通过排序 + 因果注意力的设计，在保持并行生成能力的同时支持 KV 缓存，实现了数量级的推理加速
    \item 在\textbf{无左右偏置}的离散数据任务（蛋白质、分子生成）上，扩散语言模型的优势尤为显著
\end{enumerate}

\subsection{工业应用现状}

截至本讲座录制时（2025 年秋季），MDLM 框架已经在工业界得到广泛应用：

\begin{itemize}
    \item \textbf{ByteDance Seed Diffusion}：基于 MDLM 框架构建
    \item \textbf{Google Gemini Diffusion}：采用扩散方法进行文本生成
    \item \textbf{Inception Labs Mercury Coder}：面向代码生成的扩散模型
    \item \textbf{NVIDIA}：使用 MDLM 进行分子生成（药物发现领域）
\end{itemize}

\subsection{拓展阅读}

\begin{itemize}
    \item Sahoo et al., ``Simple and Effective Masked Diffusion Language Models'' (MDLM), NeurIPS 2024
    \item Sahoo et al., ``Esoteric Language Models'' (SoLM), 2025
    \item ``Diffusion Duality'': 蒸馏算法，10 步生成 1000 词
    \item LaDa: Large Language Diffusion Models（8B 规模 MDLM 验证）
    \item BD3LM: Block Diffusion Language Models
    \item GenMol: 基于 MDLM 的分子生成
    \item KAIST CS492D 课程主页：\url{https://mhsung.github.io/kaist-cs492d-fall-2025/}
\end{itemize}

%% --- Fin del contenido --- %%

\end{document}
